\chapter{Solución de problemas}\label{cap:problemas}

Este capítulo reúne los fallos más frecuentes al instalar, operar y mantener
Qadra. Cada fila indica el síntoma observable, su causa probable y la acción
recomendada, con la ruta del documento que contiene el procedimiento completo.
Los mensajes que ve el usuario final en pantalla no se repiten aquí: están en
el capítulo \enquote{Mensajes del sistema y solución de problemas} del Manual
de usuario (parte \enquote{Referencia rápida}).

\section{Cómo diagnosticar}\label{sec:prob-diagnostico}

Toda respuesta de error de la API usa el mismo sobre JSON, con un código
estable y un mensaje genérico que no expone detalles internos ni secretos:

\begin{verbatim}
{"error":{"code":"permission_denied","message":"permission denied"}}
\end{verbatim}

El código es el dato que debe buscarse en \archivo{docs/http-api.md} (sección
de errores) y en las tablas siguientes. En una instalación desplegada, el
estado de los servicios y los diagnósticos del servidor se consultan con los
comandos de la \cref{sec:op-estado}.

\nota{Un mensaje de error no confirma que una escritura anterior haya fallado
ni que haya tenido éxito. Ante una respuesta perdida, consulte el estado
actual del registro antes de repetir un alta, un sellado o una confirmación.}

\needspace{9\baselineskip}
\section{Arranque e instalación}\label{sec:prob-arranque}

{\small
\begin{longtable}{L{3.7cm} L{4.6cm} L{7.1cm}}
  \caption{Problemas de arranque e instalación.}\label{tab:prob-arranque} \\
  \toprule
  \textbf{Síntoma} & \textbf{Causa probable} & \textbf{Acción} \\
  \midrule
  \endfirsthead
  \toprule
  \textbf{Síntoma} & \textbf{Causa probable} & \textbf{Acción} \\
  \midrule
  \endhead
  \bottomrule
  \endfoot
    \texttt{serve} no arranca o rechaza la conexión a la base
      & El esquema no está aplicado, o \texttt{DATABASE\_URL} apunta a un rol
        que puede administrar o reescribir la auditoría. El servidor no
        ejecuta DDL.
      & Aplique \texttt{database migrate -{}-runtime-role} con una conexión
        administrativa y después cambie \texttt{DATABASE\_URL} al rol
        de ejecución (\cref{cap:instalacion};
        \archivo{docs/database-operations.md}). \\
    Un guion de \archivo{pki/} termina con \texttt{CA not initialized} o
    \texttt{certificate already exists}
      & \texttt{PKI\_CA\_DIR} o \texttt{TSA\_DIR} no están exportadas en
        esa terminal, o el certificado ya se emitió.
      & Exporte ambas variables (\cref{sec:inst-pki}) y use el certificado
        existente. \\
    \texttt{npm run dev} se detiene con \texttt{artifact is stale} o
    por un archivo ausente en \archivo{src/generated/}
      & Un módulo WASM no se ha generado o no corresponde a sus fuentes.
      & Ejecute los dos guiones \texttt{scripts/build-*.sh} de la
        \cref{sec:inst-web}. \\
    El alta inicial responde \texttt{409 bootstrap\_closed}
      & La base ya contiene un usuario.
      & Inicie sesión con la cuenta existente (\cref{sec:inst-owner}). \\
    El arranque falla al comprobar la biblioteca qpdf
      & Falta \texttt{DOCUMENT\_QPDF\_LIBRARY} (o \texttt{-{}-qpdf-library}),
        o no apunta a la biblioteca qpdf 12.4.1 fijada. El validador nativo
        requiere Linux x86\_64.
      & Prepare la biblioteca con
        \archivo{scripts/setup-document-formats.sh} y exporte su salida en
        esa variable (\archivo{docs/document-format-operations.md}). \\
    El arranque o la carga de audio y video falla por el decodificador
      & Falta la instalación fijada bajo \archivo{/opt/tt-media}; un
        \texttt{ffmpeg} del sistema no la sustituye.
      & Ejecute la comprobación mostrada bajo esta tabla. Si usa otro prefijo,
        indíquelo con \texttt{TT\_FFPROBE\_PATH} y \texttt{TT\_FFMPEG\_PATH}
        (\archivo{docs/media-decoder-setup.md}). \\
    Los documentos existentes no se descifran después de reinstalar
      & Se generó una \texttt{KEK\_BASE64} nueva. Una clave nueva no descifra
        los datos existentes.
      & Restaure la KEK original de esa instalación. No genere otra para
        conseguir el arranque (\cref{cap:configuracion}). \\
\end{longtable}}

\begin{verbatim}
python3 -B scripts/install_media_decoder.py --verify --prefix /opt/tt-media
\end{verbatim}

\section{Acceso y segundo factor}\label{sec:prob-acceso}

Los rechazos de autenticación responden siempre \texttt{401} con
\texttt{mfa\_rejected}, sin motivo en el cuerpo. El motivo interno queda en el
registro del servicio, en el evento \texttt{MFA access attempt}, y se
correlaciona con la cabecera \texttt{X-Request-Id} de la respuesta; el
comando para filtrarlo está en la \cref{sec:op-estado}.

{\small
\begin{longtable}{L{3.7cm} L{4.6cm} L{7.1cm}}
  \caption{Problemas de acceso y MFA.}\label{tab:prob-acceso} \\
  \toprule
  \textbf{Síntoma} & \textbf{Causa probable} & \textbf{Acción} \\
  \midrule
  \endfirsthead
  \toprule
  \textbf{Síntoma} & \textbf{Causa probable} & \textbf{Acción} \\
  \midrule
  \endhead
  \bottomrule
  \endfoot
    \texttt{401 mfa\_rejected} con motivo
      \archivo{invalid_code_or_outside_window} o
      \archivo{code_already_used}
      & Código TOTP incorrecto, fuera de ventana o ya utilizado. El
        verificador no distingue las dos primeras causas.
      & Pida al usuario iniciar sesión otra vez con un código nuevo y revisar
        la hora de su dispositivo. No se amplían las ventanas TOTP. \\
    Motivo \archivo{challenge_expired_consumed_or_unknown}
      & El desafío venció, ya se consumió o no existe. El desafío se consume
        antes de comprobar el segundo factor.
      & El usuario debe empezar de nuevo desde correo y contraseña. \\
    Motivo \archivo{credentials_changed} o \archivo{account_inactive}
      & La cuenta cambió entre los dos pasos (incluye cambios de rol) o está
        inactiva.
      & Revise el estado y el rol de la cuenta en \textbf{Equipo} y solicite
        un nuevo inicio de sesión. \\
    No aparecen eventos \texttt{MFA access attempt} en el registro
      & El filtro de \texttt{RUST\_LOG} suprime el destino
        \texttt{qadra::access}.
      & Conserve ese destino a nivel \texttt{info}, por ejemplo
        \texttt{RUST\_LOG=info,qadra::access=info}
        (\archivo{docs/mfa-access-logs.md}). \\
    \texttt{429 account\_locked} al iniciar sesión
      & Ventana de inicio de sesión bloqueada por intentos fallidos.
      & Espere a que termine la ventana. Al restaurar Redis no borre las
        claves \archivo{identity:password-failures:*}. \\
    La recuperación de contraseña responde \texttt{503}
      & Los componentes de recuperación y entrega de correo no están
        configurados en ese servidor.
      & Revise \texttt{TT\_PASSWORD\_RESET\_ENABLED} y la configuración de
        correo (\archivo{docs/password-reset-public.md}). \\
    Un alta de integrante quedó sin confirmar al vencer la sesión
      & Se perdió la respuesta del envío. Autenticarse de nuevo no crea la
        cuenta ni reenvía la solicitud.
      & En \textbf{Equipo}, use \textbf{Retomar alta de integrante} y
        \textbf{Consultar directorio para esta cuenta} antes de decidir
        (\archivo{docs/member-enrollment-recovery.md}). Los demás
        borradores se recuperan solo si la pestaña no se recargó
        (\archivo{docs/session-editor-recovery.md}). \\
\end{longtable}}

\needspace{22\baselineskip}
\section{Documentos y firma}\label{sec:prob-documentos}

{\small
\begin{longtable}{L{3.7cm} L{4.6cm} L{7.1cm}}
  \caption{Problemas con documentos, sellado y firma.}\label{tab:prob-documentos} \\
  \toprule
  \textbf{Síntoma} & \textbf{Causa probable} & \textbf{Acción} \\
  \midrule
  \endfirsthead
  \toprule
  \textbf{Síntoma} & \textbf{Causa probable} & \textbf{Acción} \\
  \midrule
  \endhead
  \bottomrule
  \endfoot
    \texttt{503} \archivo{signing_trust_unavailable} al vincular un
      certificado o firmar
      & No hay una revisión de confianza publicada (raíz y CRL) que el
        servidor pueda usar. No se escribió ningún estado.
      & Compruebe la publicación en el portal de la CA
        (\textbf{Publicación}, \textbf{Actualizar estado}) y publique la CRL
        (\cref{cap:ca}). Después repita la operación. \\
    \texttt{409} \archivo{signing_trust_changed}
      & La revisión de confianza cambió entre la preparación y la
        confirmación, por ejemplo por una renovación de CRL.
      & Reinicie la firma desde su preparación. No reenvíe la confirmación
        anterior (\archivo{docs/http-individual-signing.md}). \\
    El sellado responde \texttt{500} y no guarda evidencia
      & La CRL cargada al iniciar el servicio superó su \texttt{nextUpdate} y
        el certificado firmante ya no se puede validar.
      & Renueve la CRL con el procedimiento de
        \archivo{docs/deployment-crl.md} y reinicie el servicio con la CRL
        renovada (\cref{cap:operacion}). \\
\end{longtable}}

\needspace{9\baselineskip}
\section{Despliegue, CRL y autoridad certificadora}\label{sec:prob-despliegue}

{\small
\begin{longtable}{L{3.7cm} L{4.6cm} L{7.1cm}}
  \caption{Problemas de despliegue, migraciones, CRL y servicio de la CA.}\label{tab:prob-despliegue} \\
  \toprule
  \textbf{Síntoma} & \textbf{Causa probable} & \textbf{Acción} \\
  \midrule
  \endfirsthead
  \toprule
  \textbf{Síntoma} & \textbf{Causa probable} & \textbf{Acción} \\
  \midrule
  \endhead
  \bottomrule
  \endfoot
    El despliegue se detiene antes de apagar o migrar
      & Cambió la huella de migraciones. Los cambios de esquema no forman
        parte del despliegue automático.
      & Realice el mantenimiento explícito: ensayo en copia aislada,
        respaldo, migración y registro de la huella
        (\archivo{docs/database-operations.md}). No edite la huella para
        eludir el rechazo. \\
    La activación de una versión falla
      & La versión nueva no superó la comprobación de salud.
      & El controlador restaura la versión anterior y devuelve el fallo; si
        tampoco la recupera, detiene la entrada web. Revise el estado y el
        registro del servicio antes de reintentar (\cref{cap:despliegue}). \\
    La API no arranca y existe \archivo{config/crl-maintenance.json}
      & Hay una renovación de CRL pendiente. El marcador bloquea arranque,
        proxy, activación y reversión.
      & No retire el marcador a mano. Revise \texttt{status} de
        \texttt{renew\_crl.py} y ejecute \texttt{resume} con el UUID de la
        operación (\archivo{docs/deployment-crl.md}). \\
    La unidad \texttt{qadra-crl-renew} queda fallida o llega una alerta
      & El recibo en \archivo{maintenance/crl-schedule/} indica
        \archivo{blocked}, \archivo{busy_near_expiry}, \archivo{ca_expiring},
        \texttt{failed} o \texttt{interrupted}.
      & Lea el recibo y actúe según su resultado. El temporizador nunca reanuda una
        operación pendiente: la recuperación es manual. \\
    Una emisión, revocación o publicación quedó con respuesta incierta
      & Se perdió la respuesta. La operación pendiente bloquea otras
        mutaciones.
      & En \textbf{Historial}, use \textbf{Retomar}: consulta o completa la
        misma operación sin cambiar sus parámetros. No inicie otra para
        eludirla (\archivo{docs/ca/operations.md}). \\
    \textbf{Publicación} muestra mantenimiento pendiente y
      \textbf{Historial} no tiene una publicación pendiente
      & El marcador lo dejó la renovación programada o una renovación por
        línea de comandos. El portal no puede retomarlo.
      & Un administrador del servidor revisa \texttt{status} de
        \texttt{renew\_crl.py} y ejecuta \texttt{resume} con el UUID
        registrado. \\
    El operador de la CA no puede entrar después de una restauración
      & La KEK del servicio no es la que protege el TOTP restaurado.
      & Restaure la configuración y la KEK del mismo conjunto. No ejecute
        \texttt{enroll-operator} ni genere otro TOTP
        (\archivo{docs/ca/restore.md}). \\
\end{longtable}}

Durante un mantenimiento de CRL no ejecute directamente
\archivo{pki/gen-crl.sh}, \archivo{pki/revoke.sh} ni la emisión de
certificados sobre esa CA: esos guiones no coordinan la publicación en la base
ni el marcador. Ante una discrepancia o un estado parcial, el bloqueo se
conserva y requiere investigación; no se bajan contadores ni se restaura SQL
para deshacer una publicación confirmada.

\needspace{9\baselineskip}
\section{Respaldo y restauración}\label{sec:prob-respaldo}

{\small
\begin{longtable}{L{3.7cm} L{4.6cm} L{7.1cm}}
  \caption{Problemas de respaldo y restauración.}\label{tab:prob-respaldo} \\
  \toprule
  \textbf{Síntoma} & \textbf{Causa probable} & \textbf{Acción} \\
  \midrule
  \endfirsthead
  \toprule
  \textbf{Síntoma} & \textbf{Causa probable} & \textbf{Acción} \\
  \midrule
  \endhead
  \bottomrule
  \endfoot
    La restauración SQL falla por roles inexistentes
      & Un \texttt{pg\_dump} de una base no incluye la creación de los roles
        del clúster.
      & Prepare los roles, propietarios y privilegios antes de restaurar
        (\archivo{docs/deployment-backups.md}). \\
    Redis ignora el RDB restaurado
      & Con AOF y RDB presentes, Redis reconstruye desde el AOF.
      & Cargue el RDB verificado en un directorio nuevo con
        \texttt{appendonly no} y habilite AOF después con
        \texttt{CONFIG SET appendonly yes}
        (\archivo{docs/deployment-backups.md}). \\
    Las sesiones anteriores siguen aceptándose tras restaurar
      & No se invalidaron sesiones ni desafíos en Redis.
      & Elimine solo las claves de sesiones y desafíos que indica la
        \cref{sec:op-redis}. No use \texttt{FLUSHDB} ni
        \texttt{FLUSHALL}; conserve los controles TOTP y de intentos. \\
    La captura de respaldo se rechaza
      & Existe una operación de la CA pendiente.
      & Termine o concilie la operación original antes de capturar. No borre
        su diario (\archivo{docs/ca/restore.md}). \\
\end{longtable}}

\importante{No desactive guardas, no regenere claves y no modifique la huella
de migraciones para conseguir que el servicio arranque después de una
restauración. Conserve intactos los datos originales hasta aceptar el
resultado (\cref{cap:operacion}).}

Si el problema no aparece en este capítulo, localice el documento del tema en
el \cref{cap:mapa} y, al solicitar soporte, comunique la operación, el código de
error, el momento y el \texttt{X-Request-Id} cuando exista. No incluya
contraseñas, códigos ni claves privadas en el reporte.
